\section{Sources and relation to earlier results}

Li proved the full-sequence criterion for Dedekind zeta functions.
Bombieri and Lagarias formulated a criterion for multisets of complex numbers
and used simultaneous phase approximation for the converse. Voros described
the exponentially oscillating contribution to the zeta coefficients when a
zero lies off the line. Lagarias developed the coefficients for automorphic
$L$-functions, including the real-part criterion used for complex characters.
These works supply the analytic setting and the mechanism by which an outer
Cayley radius affects the coefficients \cite{li,bl,voros,lag}.

The closest result in ProveIt takes dilated blocks
$\{M_j,2M_j,\ldots,jM_j\}$ for Dedekind zeta functions. It proves a positive
proportion of negative samples in each sufficiently late block if RH fails.
It also obtains the limiting normalized minimum $-m$, with $m$ the total
multiplicity on the outermost circle. The same packet uses second differences
of Li coefficients to distinguish arithmetically inequivalent fields.
Two later packets concern identification of a primitive character from
filtered fixed-stride Li samples, and simultaneous recurrence of translated
smoothed prime errors \cite{prior}.
The present classification gives a necessary and sufficient condition on an
arbitrary sampling set and realizes failure of that condition by a finite
zero multiset. For synchronized polynomial schedules the obstruction becomes
a congruence condition, with a unit constraint for prime parameters.

Bohr neighborhoods and their description by trigonometric polynomials are
classical \cite{bfw}. Weyl's theorem and Rhin's prime analogue supply the
equidistribution inputs \cite{weyl,rhin}. Bergelson, Leibman and Lesigne
proved multiple recurrence for jointly intersective polynomials.
L\^e and Spencer studied quantitative simultaneous Diophantine approximation
for intersective systems and prime analogues \cite{bll,ls}. The contribution
claimed here is the characterization for Li sums, including its converse
within the class of zero multisets. The comparison covered the 61 historical
packet READMEs and 408 TeX files in the research directories, together with
55 TeX files in 39 incoming archives at the pinned ProveIt revision.
A targeted literature search located the preceding analytic and recurrence
results. Priority for the combined theorem remains open to further comparison.

Johansson had already computed rigorous Keiper coefficients through index
100,000 \cite{joh}. For a positive index $n$, multiplying his coefficient by $n$
gives the Li normalization used here. Second differencing the published
decimals gives $d_{117}\approx0.000200147$ and
$d_{118}\approx-0.000271044$, independently confirming the initial
change of sign. Our contribution concerning this finite sign pattern is its
simultaneous recurrence at polynomially separated indices and a certificate
using forty conjugate pairs and residual moments. The computations at the
four squared indices in Table~\ref{tab:seeds} extend far beyond this table.

The survey article connects the critical-line formulation with Li
positivity \cite{wiki}. In the compact Lean edition, \texttt{src/Statements.lean}
specifies the Dirichlet zero theorem and \texttt{src/Main.lean} exports
\texttt{LFunction\_eq\_zero\_re}. The proof obtains a one-sided bound for a
weighted prime sum and passes to nonvanishing in the right half-plane.
Here its critical-line conclusion places the Cayley transforms of the
Dirichlet zeros on the unit circle. Sutherland records kernel replay and
independent checks of the exported statements for both editions \cite{lean}.
Those checks concern the upstream theorems. The sampling results above have
the mathematical proofs given in this paper.

\begin{thebibliography}{99}
\setlength{\itemsep}{2pt}
\bibitem{li} X.-J. Li, The positivity of a sequence of numbers and the
Riemann hypothesis, \emph{J. Number Theory} 65 (1997), 325--333.
\href{https://doi.org/10.1006/jnth.1997.2137}{Publisher}.
\bibitem{bl} E. Bombieri and J. C. Lagarias, Complements to Li's criterion
for the Riemann hypothesis, \emph{J. Number Theory} 77 (1999), 274--287.
\href{https://websites.umich.edu/~lagarias/doc/bombieri.pdf}{Author's copy}.
\bibitem{voros} A. Voros, Sharpenings of Li's criterion for the Riemann
hypothesis, \emph{Math. Phys. Anal. Geom.} 9 (2006), 53--63.
\href{https://doi.org/10.1007/s11040-005-9002-8}{Publisher}.
\bibitem{lag} J. C. Lagarias, Li coefficients for automorphic $L$-functions,
\emph{Ann. Inst. Fourier} 57 (2007), 1689--1740.
\href{https://doi.org/10.5802/aif.2311}{Publisher}.
\bibitem{bfw} V. Bergelson, H. Furstenberg and B. Weiss, Piecewise-Bohr sets
of integers and combinatorial number theory, Definition~1.1.
\href{https://people.math.osu.edu/bergelson.1/BFW_Bohr.pdf}{Author's copy}.
\bibitem{weyl} H. Weyl, \"Uber die Gleichverteilung von Zahlen mod. Eins,
\emph{Math. Ann.} 77 (1916), 313--352.
\href{https://doi.org/10.1007/BF01475864}{Publisher}.
\bibitem{rhin} G. Rhin, Sur la r\'epartition modulo $1$ des suites $f(p)$,
\emph{Acta Arith.} 23 (1973), 217--248.
\href{https://doi.org/10.4064/aa-23-3-217-248}{Publisher}.
\bibitem{bll} V. Bergelson, A. Leibman and E. Lesigne, Intersective
polynomials and the polynomial Szemer\'edi theorem,
\emph{Adv. Math.} 219 (2008), 369--388.
\href{https://arxiv.org/abs/0710.4862}{Preprint}.
\bibitem{ls} T. H. L\^e and C. V. Spencer, Intersective polynomials and
Diophantine approximation, \emph{Int. Math. Res. Not.} 2014 (2014), 1153--1173.
\href{https://doi.org/10.1093/imrn/rns242}{Publisher}.
\bibitem{joh} F. Johansson, Rigorous high-precision computation of the Hurwitz
zeta function and its derivatives, \emph{Numer. Algorithms} 69 (2015), 253--270.
\href{https://fredrikj.net/math/hurwitz_zeta.html}{Paper and coefficient data}.
\bibitem{prior} V. Reshetnikov, \emph{ProveIt}, Generalized Riemann Hypothesis
research collection, revision \texttt{4d7990202f9f}.
\href{https://github.com/VladimirReshetnikov/ProveIt/tree/4d7990202f9f106b2ffe8c209f073ad01d5e1987/Analysis/GeneralizedRiemannHypothesis}{Collection}.
\href{https://github.com/VladimirReshetnikov/ProveIt/blob/4d7990202f9f106b2ffe8c209f073ad01d5e1987/Analysis/GeneralizedRiemannHypothesis/Research/grh_method_transfers_research/li.tex}{Sparse dilation-block criterion}.
\href{https://github.com/VladimirReshetnikov/ProveIt/blob/4d7990202f9f106b2ffe8c209f073ad01d5e1987/Analysis/GeneralizedRiemannHypothesis/Research/grh_method_transfers_research/rigidity.tex}{Arithmetic equivalence}.
\href{https://github.com/VladimirReshetnikov/ProveIt/blob/4d7990202f9f106b2ffe8c209f073ad01d5e1987/Analysis/GeneralizedRiemannHypothesis/docs/incoming/dirichlet_li_rigidity.zip}{Filtered Li samples}.
\href{https://github.com/VladimirReshetnikov/ProveIt/blob/4d7990202f9f106b2ffe8c209f073ad01d5e1987/Analysis/GeneralizedRiemannHypothesis/docs/incoming/prime_peak_rigidity.zip}{Smoothed prime errors}.
\bibitem{lean} L. Sernau and collaborators, shorter editions of the
Alp\"oge--Claude proof of RH and GRH for Dirichlet $L$-functions,
revision \texttt{72fead102f66} (2026).
\href{https://github.com/lsernau/grh-lean/tree/72fead102f663d5c7a407964a13dc0cf9d977207}{Source}.
A. V. Sutherland, \href{https://math.mit.edu/~drew/shortgrh/}{independent verification records}.
\bibitem{wiki} Wikipedia contributors, \emph{Riemann hypothesis}, consulted
26 November 2026. \href{https://en.wikipedia.org/wiki/Riemann_hypothesis}{Article}.
\end{thebibliography}
